\documentclass[11pt, a4paper]{article}

% ---------- Codifica e lingua ----------
\usepackage[utf8]{inputenc}
\usepackage[T1]{fontenc}
\usepackage[italian]{babel}

% ---------- Matematica e font ----------
\usepackage{amsmath, amssymb, amsthm}
\usepackage{mathtools}
\usepackage{newpxtext, newpxmath}
\usepackage[scale=0.95]{sourcesans}
\usepackage[scale=0.95]{inconsolata}
\usepackage{bm}
\usepackage{microtype}
\linespread{1.08}

% ---------- Pagina e colori ----------
\usepackage[left=2.3cm, right=2.3cm, top=2.3cm, bottom=2.3cm]{geometry}
\usepackage[table]{xcolor}

% Palette: nero-verde di base, verde come accento, arancione solo per ciò che è cruciale
\definecolor{sfondo}{HTML}{0D110E}      % nero con una punta di verde
\definecolor{sfondobox}{HTML}{151C17}   % fondo dei riquadri
\definecolor{bordo}{HTML}{2B3B30}       % linee e filetti
\definecolor{testo}{HTML}{C8CEC9}       % grigio chiaro (non bianco)
\definecolor{tenue}{HTML}{87928A}       % grigio secondario
\definecolor{verde}{HTML}{5CC98A}       % accento principale
\definecolor{verdechiaro}{HTML}{9FDDB4} % accento secondario (esempi)
\definecolor{verdeacido}{HTML}{B4D46A}  % proprietà
\definecolor{verdescuro}{HTML}{3C8F5F}  % dettagli
\definecolor{arancio}{HTML}{F0954A}     % solo per l'essenziale

\pagecolor{sfondo}
\color{testo}

% ---------- Altri pacchetti ----------
\usepackage{parskip}
\usepackage{graphicx}
\usepackage{forest}
\usetikzlibrary{mindmap, trees, arrows.meta, automata, positioning, shapes}
\usepackage{listings}
\usepackage{enumitem}
\usepackage{booktabs}
\usepackage{array}
\usepackage{titlesec}
\usepackage[most]{tcolorbox}
\usepackage[hidelinks]{hyperref}

\forestset{
  default preamble={
    for tree={
      draw=bordo,
      fill=sfondobox,
      rounded corners,
      node options={align=center},
      font=\sffamily\small,
      edge={color=verdescuro, line width=1pt}
    }
  }
}

% ---------- Titoli ----------
% Il titolo di sezione è arancione (la cosa più importante della pagina),
% la numerazione e i sottotitoli sono verdi.
\renewcommand{\thesubsection}{\arabic{subsection}}

\titleformat{\section}
  {\sffamily\bfseries\LARGE\color{arancio}}{}{0pt}{}
  [\vspace{2pt}{\color{verdescuro}\titlerule[1.2pt]}]
\titleformat{\subsection}
  {\sffamily\bfseries\Large\color{testo}}{\textcolor{verde}{\thesubsection}}{0.6em}{}
\titleformat{\subsubsection}
  {\sffamily\bfseries\large\color{verde}}{}{0pt}{}

\titlespacing*{\section}{0pt}{0.5\baselineskip}{0.8\baselineskip}
\titlespacing*{\subsection}{0pt}{1.5\baselineskip}{0.4\baselineskip}
\titlespacing*{\subsubsection}{0pt}{1\baselineskip}{0.2\baselineskip}

% ---------- Elenchi ----------
\setlist{topsep=2pt, itemsep=2pt, parsep=0pt}
\setlist[itemize,1]{label=\textcolor{verde}{\textbullet}, leftmargin=1.3em}
\setlist[itemize,2]{label=\textcolor{tenue}{--}, leftmargin=1.2em}
\setlist[enumerate,1]{label=\textcolor{verde}{\sffamily\bfseries\arabic*}, leftmargin=1.3em}

% ---------- Testo ----------
\renewcommand{\emph}[1]{\textcolor{verdechiaro}{\itshape #1}}

% ---------- Box ----------
\tcbset{
  appunti/.style={
    enhanced, breakable=false,
    colback=sfondobox, coltext=testo,
    boxrule=0pt, frame hidden, arc=2pt,
    left=10pt, right=10pt, top=7pt, bottom=7pt,
    before skip=9pt, after skip=11pt,
    fonttitle=\sffamily\bfseries,
    detach title,
    before upper={\setlength{\parskip}{4pt}\tcbtitle\par},
    borderline west={2.5pt}{0pt}{#1},
    coltitle=#1
  }
}

\newtcolorbox{defbox}[1]{appunti=verde, title={#1}}
\newtcolorbox{esempio}[1][Esempio]{appunti=verdechiaro, title={#1}}
\newtcolorbox{riepilogo}[1]{appunti=verdescuro, title={#1}}
\newtcolorbox{proprietà}[1]{appunti=verdeacido, title={#1}}
% Riservato ai risultati fondamentali (da sapere a memoria)
\newtcolorbox{chiave}[1]{appunti=arancio, title={#1},
  colback=arancio!6!sfondobox}

% ---------- Pseudocodice ----------
\lstdefinelanguage{pseudo}{
  sensitive=true,
  alsoletter={-?},
  morekeywords={function, returns, return, if, then, else, do, and, static, inputs, while, an, loop, for, each, in, unless, local, variables},
  emph={true, false},
  emphstyle=\color{verdeacido}
}

\lstdefinestyle{pseudostile}{
  language=pseudo,
  basicstyle=\ttfamily\small\color{testo},
  keywordstyle=\color{verde}\bfseries,
  columns=fullflexible,
  keepspaces=true,
  breaklines=true,
  breakindent=2.5em,
  mathescape=true,
  aboveskip=4pt, belowskip=0pt
}

\newtcblisting{codice}[1]{
  appunti=bordo, coltitle=tenue, title={#1},
  fonttitle=\sffamily\small\bfseries,
  listing only,
  listing options={style=pseudostile}
}

% ---------- Tabelle ----------
\arrayrulecolor{bordo}
\newcommand{\intest}[1]{\textcolor{verde}{\sffamily\bfseries #1}}
\newcommand{\vero}{\textcolor{verde}{V}}
\newcommand{\falso}{\textcolor{tenue}{F}}
\newcommand{\nota}[1]{\textcolor{tenue}{#1}}

% ---------- Macro ----------
\newcommand{\KB}{\mathrm{KB}}
\newcommand{\Pv}{\mathbf{P}}               % distribuzione (vettore)
\newcommand{\tup}[1]{\langle #1 \rangle}    % vettore di probabilità
\newcommand{\known}{\mathit{known}}
\newcommand{\fringe}{\mathit{fringe}}
\newcommand{\other}{\mathit{other}}
\newcommand{\unknown}{\mathit{unknown}}


\begin{document}

% ---------- Intestazione ----------
\begingroup
  \raggedright
  {\sffamily\bfseries\Huge\color{arancio} Intelligenza artificiale\par}
  \vspace{4pt}
  {\sffamily\color{tenue} Giorgio Poluzzi \quad·\quad \today\par}
\endgroup
\vspace{1.6\baselineskip}

% =====================================================================
\section{Logical Agents}
% =====================================================================

\subsection{Agente logico}

Un \textbf{agente basato sulla conoscenza} (agente logico) è costituito da due componenti distinte:

\begin{itemize}
  \item \textbf{Knowledge base (KB)}: un insieme di \emph{frasi} espresse in un linguaggio formale; rappresenta ciò che l'agente sa del dominio;
  \item \textbf{Motore inferenziale}: un insieme di algoritmi generali, indipendenti dal dominio, che derivano nuove frasi a partire da quelle già presenti nella KB.
\end{itemize}

L'approccio è \textbf{dichiarativo}: non si programma il comportamento dell'agente passo per passo, ma gli si comunica ciò che deve sapere (\textbf{Tell}) e gli si chiede che cosa fare (\textbf{Ask}); sarà l'inferenza a ricavare la risposta.

Un agente si può descrivere a due livelli:
\begin{itemize}
  \item \textbf{livello della conoscenza}: \emph{che cosa} sa l'agente;
  \item \textbf{livello implementativo}: \emph{come} tale conoscenza è realizzata (strutture dati e algoritmi).
\end{itemize}

\subsubsection*{Il ciclo dell'agente}

A ogni istante $t$ l'agente esegue il medesimo ciclo:
\begin{enumerate}
  \item \textbf{Tell}: comunica alla KB la percezione ricevuta al tempo $t$;
  \item \textbf{Ask}: interroga la KB su quale azione compiere;
  \item \textbf{Tell}: registra nella KB l'azione scelta;
  \item incrementa $t$ e restituisce l'azione.
\end{enumerate}

Affinché questo schema funzioni, l'agente deve essere in grado di:
\begin{itemize}
  \item rappresentare stati e azioni;
  \item incorporare nuove percezioni;
  \item aggiornare la propria rappresentazione del mondo;
  \item dedurre le proprietà del mondo che non osserva direttamente;
  \item dedurre quali azioni siano appropriate.
\end{itemize}

\subsection{Logica, modelli, entailment e inferenza}

\begin{defbox}{Logica}
  Linguaggio formale per rappresentare informazioni in modo che se ne possano trarre conclusioni. È definita da due elementi:
  \begin{itemize}
    \item \textbf{sintassi}: stabilisce quali frasi sono ben formate;
    \item \textbf{semantica}: attribuisce un significato alle frasi, cioè stabilisce quando una frase è vera in un dato modello.
  \end{itemize}
\end{defbox}

\begin{defbox}{Modello}
  Un mondo formalizzato rispetto al quale è possibile stabilire se una frase è vera o falsa.
  \begin{itemize}
    \item $m$ è un \textbf{modello} di $\alpha$ se $\alpha$ è vera in $m$;
    \item $M(\alpha)$ indica l'insieme di tutti i modelli di $\alpha$.
  \end{itemize}
\end{defbox}

\begin{chiave}{Entailment (conseguenza logica)}
  $\KB \models \alpha$ (``la KB implica logicamente $\alpha$'') significa che $\alpha$ è vera in \textbf{ogni} modello in cui la KB è vera. In termini insiemistici:
  \[
    \KB \models \alpha \iff M(\KB) \subseteq M(\alpha)
  \]
\end{chiave}

\begin{esempio}
  Se $\KB = \{P \land Q\}$ e $\alpha = P$, allora $\KB \models \alpha$: in ogni modello in cui $P \land Q$ è vera, anche $P$ è vera.
\end{esempio}

\begin{defbox}{Inferenza}
  $\KB \vdash_i \alpha$ significa che $\alpha$ può essere \emph{derivata} dalla KB mediante la procedura $i$.
\end{defbox}

\textbf{Perché distinguere i due concetti.} L'entailment è una relazione \emph{semantica}: dice che $\alpha$ segue dalla KB, ma non dice come trovarla. L'inferenza è un procedimento \emph{sintattico}: è il metodo concreto con cui l'agente costruisce $\alpha$ manipolando le frasi.

Una procedura di inferenza $i$ si giudica in base a due proprietà:

\begin{center}
  \renewcommand{\arraystretch}{1.35}
  \begin{tabular}{>{\bfseries}l l l}
    \toprule
    \intest{Proprietà} & \intest{Significato} & \intest{Formula} \\
    \midrule
    Correttezza & deriva \textbf{soltanto} frasi implicate & se $\KB \vdash_i \alpha$, allora $\KB \models \alpha$ \\
    Completezza & deriva \textbf{tutte} le frasi implicate & se $\KB \models \alpha$, allora $\KB \vdash_i \alpha$ \\
    \bottomrule
  \end{tabular}
\end{center}

\clearpage
\subsection{Logica proposizionale}

È la logica più semplice: il suo mondo è descritto unicamente da proposizioni vere o false.

\subsubsection*{Sintassi}

I simboli $P_1, P_2, \dots$ sono \textbf{frasi atomiche}. Date le frasi $S, S_1, S_2$, si costruiscono frasi complesse mediante i connettivi:

\begin{center}
  \renewcommand{\arraystretch}{1.25}
  \begin{tabular}{l l}
    \toprule
    \intest{Frase} & \intest{Nome} \\
    \midrule
    $\lnot S$ & negazione \\
    $S_1 \land S_2$ & congiunzione (``e'') \\
    $S_1 \lor S_2$ & disgiunzione (``o'') \\
    $S_1 \Rightarrow S_2$ & implicazione \\
    $S_1 \Leftrightarrow S_2$ & bicondizionale (``se e solo se'') \\
    \bottomrule
  \end{tabular}
\end{center}

\subsubsection*{Semantica}

Un \textbf{modello} assegna un valore di verità, vero (\vero) o falso (\falso), a ogni simbolo; ad esempio $P_{1,2}=F$, $P_{2,2}=F$, $P_{3,1}=V$. Con $n$ simboli esistono $2^n$ modelli (con 3 simboli, $2^3 = 8$), che possono essere enumerati in modo automatico.

\textbf{Regole di valutazione:}
\begin{itemize}
  \item $\lnot S$ è vera se e solo se $S$ è falsa;
  \item $S_1 \land S_2$ è vera se e solo se entrambe sono vere;
  \item $S_1 \lor S_2$ è vera se e solo se almeno una è vera;
  \item $S_1 \Rightarrow S_2$ è falsa \textbf{soltanto} quando $S_1$ è vera e $S_2$ è falsa;
  \item $S_1 \Leftrightarrow S_2$ è vera quando valgono entrambe le implicazioni $S_1 \Rightarrow S_2$ e $S_2 \Rightarrow S_1$.
\end{itemize}

\textbf{Tavola di verità:}

\begin{center}
  \renewcommand{\arraystretch}{1.25}
  \begin{tabular}{*{2}{>{\centering\arraybackslash}p{0.9cm}}@{\hspace{0.8cm}}*{5}{>{\centering\arraybackslash}p{1.7cm}}}
    \toprule
    $S_1$ & $S_2$ & $\lnot S_1$ & $S_1 \land S_2$ & $S_1 \lor S_2$ & $S_1 \Rightarrow S_2$ & $S_1 \Leftrightarrow S_2$ \\
    \midrule
    \falso & \falso & \vero & \falso & \falso & \vero & \vero \\
    \falso & \vero & \vero & \falso & \vero & \vero & \falso \\
    \vero & \falso & \falso & \falso & \vero & \falso & \falso \\
    \vero & \vero & \falso & \vero & \vero & \vero & \vero \\
    \bottomrule
  \end{tabular}
\end{center}

\begin{esempio}[Esempio di valutazione ricorsiva]
  Con il modello indicato sopra:
  \[
    \lnot P_{1,2} \land (P_{2,2} \lor P_{3,1}) = V \land (F \lor V) = V \land V = V
  \]
\end{esempio}

\clearpage
\subsection{Inferenza per enumerazione}

Per verificare $\KB \models \alpha$ si enumerano tutti i modelli mediante una \textbf{ricerca in profondità}: a ogni livello si fissa un simbolo prima a vero e poi a falso. L'entailment vale se $\alpha$ risulta vera in tutti i modelli in cui è vera la KB. L'algoritmo è corretto e completo, ma ha costo $O(2^n)$ nel tempo.

\begin{codice}{TT-Entails? (1/2)}
function TT-Entails?(KB, $\alpha$) returns true or false
  inputs: KB, $\text{knowledge base, a sentence in prop logic}$
          $\alpha$, $\text{a query, sentence in prop logic}$

  symbols $\leftarrow$ $\text{a list of the proposition symbols in KB and }\alpha$
  return TT-Check-All(KB, $\alpha$, symbols, [])
\end{codice}

\begin{codice}{TT-Check-All (2/2)}
function TT-Check-All(KB, $\alpha$, symbols, model) returns true or false
  if Empty?(symbols) then
    if PL-True?(KB, model) then return PL-True?($\alpha$, model)
    else return true
  else do
    P $\leftarrow$ First(symbols); rest $\leftarrow$ Rest(symbols)
    return TT-Check-All(KB, $\alpha$, rest, Extend(P, true, model))
       and TT-Check-All(KB, $\alpha$, rest, Extend(P, false, model))
\end{codice}

\subsection{Metodi di dimostrazione}

\begin{tcbraster}[raster columns=2, raster equal height, raster column skip=10pt, raster before skip=8pt, raster after skip=10pt]
  \begin{defbox}{Model checking}
    \raggedright
    \begin{itemize}
      \item enumerazione con tavole di verità;
      \item backtracking (DPLL);
      \item ricerca locale euristica (min-conflicts, hill-climbing).
    \end{itemize}
  \end{defbox}
  \begin{defbox}{Applicazione di regole di inferenza}
    \raggedright
    \begin{itemize}
      \item si generano nuove frasi a partire da quelle esistenti;
      \item una dimostrazione è una sequenza di applicazioni di regole;
      \item in genere richiede di portare le frasi in una forma normale.
    \end{itemize}
  \end{defbox}
\end{tcbraster}

\subsection{Equivalenza logica}

\begin{defbox}{Definizione}
  Due frasi sono logicamente equivalenti se e solo se sono vere negli stessi modelli, ossia se ciascuna implica l'altra:
  \[
    \alpha \equiv \beta \iff \alpha \models \beta \ \text{ e } \ \beta \models \alpha
  \]
\end{defbox}

\begin{center}
  \renewcommand{\arraystretch}{1.3}
  \begin{tabular}{>{$}l<{$} >{$}c<{$} >{$}l<{$} l}
    \toprule
    \multicolumn{3}{l}{\intest{Equivalenza}} & \intest{Nome} \\
    \midrule
    (\alpha \land \beta) & \equiv & (\beta \land \alpha) & commutatività di $\land$ \\
    (\alpha \lor \beta) & \equiv & (\beta \lor \alpha) & commutatività di $\lor$ \\
    \addlinespace[4pt]
    ((\alpha \land \beta) \land \gamma) & \equiv & (\alpha \land (\beta \land \gamma)) & associatività di $\land$ \\
    ((\alpha \lor \beta) \lor \gamma) & \equiv & (\alpha \lor (\beta \lor \gamma)) & associatività di $\lor$ \\
    \addlinespace[4pt]
    \lnot(\lnot\alpha) & \equiv & \alpha & eliminazione della doppia negazione \\
    \addlinespace[4pt]
    (\alpha \Rightarrow \beta) & \equiv & (\lnot\beta \Rightarrow \lnot\alpha) & contrapposizione \\
    (\alpha \Rightarrow \beta) & \equiv & (\lnot\alpha \lor \beta) & eliminazione dell'implicazione \\
    (\alpha \Leftrightarrow \beta) & \equiv & ((\alpha \Rightarrow \beta) \land (\beta \Rightarrow \alpha)) & eliminazione del bicondizionale \\
    \addlinespace[4pt]
    \lnot(\alpha \land \beta) & \equiv & (\lnot\alpha \lor \lnot\beta) & De Morgan \\
    \lnot(\alpha \lor \beta) & \equiv & (\lnot\alpha \land \lnot\beta) & De Morgan \\
    \addlinespace[4pt]
    (\alpha \land (\beta \lor \gamma)) & \equiv & ((\alpha \land \beta) \lor (\alpha \land \gamma)) & distributività di $\land$ su $\lor$ \\
    (\alpha \lor (\beta \land \gamma)) & \equiv & ((\alpha \lor \beta) \land (\alpha \lor \gamma)) & distributività di $\lor$ su $\land$ \\
    \bottomrule
  \end{tabular}
\end{center}

\subsection{Validità e soddisfacibilità}

\begin{tcbraster}[raster columns=3, raster equal height, raster column skip=10pt, raster before skip=8pt, raster after skip=10pt]
  \begin{defbox}{Validità}
    \raggedright
    Una frase è \textbf{valida} (tautologia) se è vera in tutti i modelli.
    \begin{esempio}
      $A \lor \lnot A$ \\
      $A \Rightarrow A$ \\
      $(A \land (A \Rightarrow B)) \Rightarrow B$
    \end{esempio}
  \end{defbox}
  \begin{defbox}{Soddisfacibilità}
    \raggedright
    Una frase è \textbf{soddisfacibile} se è vera in almeno un modello.
    \begin{esempio}
      $A \lor B$ \\
      $C$
    \end{esempio}
  \end{defbox}
  \begin{defbox}{Insoddisfacibilità}
    \raggedright
    Una frase è \textbf{insoddisfacibile} se non è vera in alcun modello.
    \begin{esempio}
      $A \land \lnot A$
    \end{esempio}
  \end{defbox}
\end{tcbraster}

\subsubsection*{Collegamento con l'inferenza}

Questi concetti consentono di ricondurre l'entailment a una verifica di validità o di insoddisfacibilità:

\begin{chiave}{Teorema di deduzione e refutazione}
  \begin{itemize}
    \item \textbf{Teorema di deduzione}: $\KB \models \alpha \iff (\KB \Rightarrow \alpha)$ è valida;
    \item \textbf{Refutazione} (dimostrazione per assurdo): $\KB \models \alpha \iff (\KB \land \lnot\alpha)$ è insoddisfacibile.
  \end{itemize}
\end{chiave}

\begin{defbox}{Sistema di inferenza}
  Un sistema di inferenza è un insieme $\mathcal{R}$ di regole della forma:
  \[
    \frac{A_1 \ \dots \ A_k}{A}
    \qquad \text{ossia} \qquad
    \frac{\text{premesse}}{\text{conclusione}}
  \]
\end{defbox}

Esiste una \textbf{derivazione} di $A$ da un insieme di formule $\Gamma$ nel sistema $\mathcal{R}$ (si scrive $\Gamma \vdash_{\mathcal{R}} A$) se esiste una successione $A_1, \dots, A_n$ tale che:
\begin{itemize}
  \item $A_n = A$;
  \item per ogni $i \in \{1, \dots, n\}$ vale una delle due condizioni:
  \begin{enumerate}
    \item $A_i \in \Gamma$, oppure
    \item $A_i$ si ottiene applicando una regola di $\mathcal{R}$ a formule che la precedono nella successione.
  \end{enumerate}
\end{itemize}

La successione $A_1, \dots, A_n$ si dice \textbf{dimostrazione} di $A$; l'insieme $\Gamma$ è l'insieme delle \emph{premesse} (o ipotesi) di $A$.

\begin{proprietà}{Proprietà dei sistemi di inferenza}
  \begin{itemize}
    \item \textbf{Correttezza}: ogni conclusione è conseguenza logica delle premesse, cioè se $\psi_1,\dots,\psi_n \vdash \psi$ allora $\psi_1,\dots,\psi_n\models\psi$.
    \item \textbf{Completezza}: se $H\models\psi$, allora esiste una derivazione di $\psi$ da $H$.
    \item \textbf{Completezza refutazionale}: se $H\cup\{\lnot\psi\}$ è insoddisfacibile, il sistema è in grado di derivarne una contraddizione.
  \end{itemize}
\end{proprietà}

\subsection{Problema della deduzione}

Si vuole stabilire se $\Gamma\models A$. Le strade principali sono due:
\begin{itemize}
  \item \textbf{refutazione}: si dimostra che $\Gamma\cup\{\lnot A\}$ è insoddisfacibile (ad esempio con la risoluzione);
  \item \textbf{forward / backward chaining}: algoritmi di costo polinomiale (anzi lineare), applicabili però soltanto a una classe ristretta di formule, le \emph{clausole di Horn}.
\end{itemize}

\begin{defbox}{Forma normale congiuntiva (CNF)}
  Una formula è in CNF se è una congiunzione ($\land$) di \textbf{clausole}, dove ogni clausola è una disgiunzione ($\lor$) di \textbf{letterali} (simboli o simboli negati):
  \[
    (\ell_{1,1} \lor \dots \lor \ell_{1,k}) \land \dots \land (\ell_{m,1} \lor \dots \lor \ell_{m,k})
  \]
  Ogni formula della logica proposizionale è convertibile in una formula equivalente in CNF.
\end{defbox}

\subsubsection*{Risoluzione}

Se una clausola contiene un letterale $\ell$ e un'altra contiene il suo opposto $\lnot\ell$, si possono eliminare entrambi e unire i letterali rimanenti in una nuova clausola (il \emph{risolvente}):
\[
  \frac{\ell_1 \lor \dots \lor \ell_k \lor \ell \qquad m_1 \lor \dots \lor m_n \lor \lnot\ell}
       {\ell_1 \lor \dots \lor \ell_k \lor m_1 \lor \dots \lor m_n}
\]
Applicata a $\KB \land \lnot\alpha$, la risoluzione è \textbf{refutazionalmente completa}: se si genera la clausola vuota (contraddizione), allora $\KB \models \alpha$.

\begin{codice}{Algoritmo di risoluzione}
function PL-Resolution(KB, $\alpha$) returns true or false
  inputs: KB, $\alpha$
  clauses $\leftarrow$ set of clauses in CNF form of KB $\wedge\lnot\alpha$
  new $\leftarrow\{\}$
  loop do
    for each C$_i$, C$_j$ in clauses do
      resolvents $\leftarrow$ PL-Resolve(C$_i$, C$_j$)
      if resolvents contains the empty clause then return true
      new $\leftarrow$ new $\cup$ resolvents
    if new $\subseteq$ clauses then return false
    clauses $\leftarrow$ clauses $\cup$ new
\end{codice}

\subsubsection*{Clausole di Horn e Modus Ponens}

Nel forward e nel backward chaining la KB è una congiunzione di \textbf{clausole di Horn}, che possono essere di due tipi:
\begin{itemize}
  \item un singolo simbolo, cioè un \textbf{fatto} (ad esempio $A$);
  \item una congiunzione di simboli che implica un simbolo, cioè una \textbf{regola} (ad esempio $C\land D\Rightarrow B$).
\end{itemize}

Su queste clausole si applica il \textbf{Modus Ponens}:
\[
  \frac{\alpha_1,\ \dots,\ \alpha_n \qquad \alpha_1\land\dots\land\alpha_n\Rightarrow\beta}{\beta}
\]
Il Modus Ponens è completo per le KB di Horn, e gli algoritmi che ne derivano hanno tempo di esecuzione lineare nella dimensione della KB.

\subsubsection*{Forward chaining}

L'idea è di attivare ogni regola le cui premesse sono tutte già note, aggiungere la sua conclusione alla KB e ripetere finché non si ottiene la query (o finché non si può dedurre nulla di nuovo).

\begin{codice}{Forward chaining}
function PL-FC-Entails?(KB, q) returns true or false
  inputs: KB, q
  local variables: count, a table indexed by clause
                   inferred, a table indexed by symbol
                   agenda, a list of symbols
  while agenda is not empty do
    p $\leftarrow$ Pop(agenda)
    unless inferred[p] do
      inferred[p] $\leftarrow$ true
      for each Horn clause c in whose premise p appears do
        decrement count[c]
        if count[c] = 0 then do
          if Head[c] = q then return true
          Push(Head[c], agenda)
  return false
\end{codice}

\textbf{Completezza.} Il forward chaining deriva ogni simbolo atomico implicato dalla KB. Infatti:
\begin{enumerate}
  \item l'algoritmo termina in un \emph{punto fisso}, in cui non è più possibile dedurre nulla di nuovo;
  \item si consideri lo stato finale come un modello $m$: vero per i simboli dedotti, falso per tutti gli altri;
  \item ogni clausola della KB è vera in $m$: se una regola $a_1\land\dots\land a_k\Rightarrow b$ fosse falsa in $m$, le premesse sarebbero tutte vere e $b$ falso; ma allora l'algoritmo potrebbe ancora dedurre $b$, contro l'ipotesi di punto fisso (contraddizione);
  \item dunque $m$ è un modello della KB; se $\KB\models a$, allora $a$ è vero in ogni modello della KB, in particolare in $m$, e quindi è stato dedotto.
\end{enumerate}

\subsubsection*{Backward chaining}

Si lavora all'indietro a partire dalla query $q$: se $q$ è già nota non c'è nulla da fare; altrimenti si cerca una regola che abbia $q$ come conclusione e si tenta di dimostrare ricorsivamente tutte le sue premesse.

Accorgimenti fondamentali:
\begin{itemize}
  \item \textbf{evitare i cicli}: controllare se un nuovo sotto-obiettivo è già presente nella pila degli obiettivi;
  \item \textbf{evitare il lavoro ripetuto}: ricordare quali sotto-obiettivi sono già stati dimostrati e quali sono falliti.
\end{itemize}

\subsubsection*{Confronto}

\begin{center}
  \renewcommand{\arraystretch}{1.35}
  \begin{tabular}{>{\bfseries}l p{5.6cm} p{5.6cm}}
    \toprule
     & \intest{Forward chaining} & \intest{Backward chaining} \\
    \midrule
    Guidato da & i dati & l'obiettivo (la query) \\
    Lavoro & deduce tutto, anche ciò che è irrilevante per la query & esplora solo ciò che è rilevante \\
    Costo & lineare nella dimensione della KB & spesso molto inferiore al lineare \\
    Adatto a & percezione, riconoscimento & problem solving, domande \\
    \bottomrule
  \end{tabular}
\end{center}

\clearpage
% =====================================================================
\section{Rappresentare l'incertezza}
% =====================================================================

\subsection{Il problema dell'incertezza}

Un agente reale raramente dispone di informazioni certe. Le cause principali sono quattro:
\begin{enumerate}
  \item \textbf{osservabilità parziale}: l'agente non vede tutto lo stato del mondo;
  \item \textbf{sensori rumorosi}: anche ciò che vede può essere misurato in modo impreciso;
  \item \textbf{esiti incerti delle azioni}: un'azione può non produrre l'effetto atteso;
  \item \textbf{complessità}: modellare e prevedere il mondo in ogni dettaglio è impraticabile.
\end{enumerate}

Di fronte a questi problemi, un approccio \textbf{puramente logico} fallisce in uno di due modi:
\begin{itemize}
  \item o \textbf{afferma il falso}, trattando come certa una conclusione che non lo è;
  \item oppure produce conclusioni talmente piene di condizioni (``\dots{} a meno che \dots{} e a meno che \dots'') da risultare \textbf{troppo deboli} per prendere una decisione.
\end{itemize}

\subsubsection*{Metodi per gestire l'incertezza}

\begin{tcbraster}[raster columns=3, raster equal height, raster column skip=10pt, raster before skip=8pt, raster after skip=10pt]
  \begin{defbox}{Logica di default}
    \raggedright
    (Logica non monotona.) Si assume che le cose vadano ``normalmente'' finché un'evidenza non lo contraddice.
    \textbf{Problemi}: quali assunzioni sono ragionevoli? Come si gestiscono le contraddizioni?
  \end{defbox}
  \begin{defbox}{Logica fuzzy}
    \raggedright
    Rappresenta \textbf{gradi di verità} (una frase è ``vera al 20\%''), \textbf{non} l'incertezza: è un concetto diverso.
  \end{defbox}
  \begin{defbox}{Probabilità}
    \raggedright
    Date le evidenze disponibili, a una proposizione si associa un \textbf{grado di credenza} numerico compreso tra 0 e 1.
  \end{defbox}
\end{tcbraster}

\subsection{La probabilità come grado di credenza}

Le asserzioni probabilistiche \emph{riassumono} due tipi di limite:
\begin{itemize}
  \item \textbf{pigrizia}: non si vogliono (o non si possono) elencare tutte le eccezioni e le condizioni;
  \item \textbf{ignoranza}: mancano fatti rilevanti, condizioni iniziali, conoscenza teorica.
\end{itemize}

Si adotta la visione \textbf{soggettiva} o \textbf{bayesiana}: la probabilità mette in relazione una proposizione con lo \emph{stato di conoscenza} dell'agente. Ne seguono due osservazioni importanti:
\begin{itemize}
  \item la probabilità esprime un \textbf{grado di credenza}, non un grado di verità: il fatto in sé è vero o falso, è l'agente a non saperlo;
  \item la probabilità di una proposizione \textbf{cambia con nuove evidenze}, perché cambia ciò che l'agente sa.
\end{itemize}

È un'analogia precisa con la logica: dire ``$P(\alpha \mid \text{evidenza}) = p$'' è simile a dire $\KB \models \alpha$, cioè descrive il rapporto tra $\alpha$ e ciò che si sa, non la verità di $\alpha$ nel mondo.

\subsection{Decidere in condizioni di incertezza}

La sola probabilità non basta per scegliere un'azione: occorre anche sapere quanto l'agente \emph{preferisce} ciascun esito. Le preferenze si rappresentano con la \textbf{teoria dell'utilità}, che assegna a ogni esito un valore numerico $U$.

\begin{chiave}{Teoria delle decisioni e principio MEU}
  \[
    \text{Teoria delle decisioni} \;=\; \text{teoria dell'utilità} \;+\; \text{teoria della probabilità}
  \]
  \textbf{Principio della massima utilità attesa (MEU)}: un agente razionale sceglie l'azione con l'utilità attesa più alta, cioè la media delle utilità degli esiti pesata sulle loro probabilità:
  \[
    a^* = \arg\max_{a} \; \sum_{s} P(s \mid a)\, U(s)
  \]
\end{chiave}

\subsection{Fondamenti di probabilità}

\begin{defbox}{Spazio campionario e modello probabilistico}
  \begin{itemize}
    \item Lo \textbf{spazio campionario} $\Omega$ è l'insieme di tutti i possibili esiti.
    \item Ogni $\omega \in \Omega$ è un \textbf{punto campionario} (o \emph{mondo possibile}, o \emph{evento atomico}).
    \item Un \textbf{modello probabilistico} è uno spazio campionario con un'assegnazione $P(\omega)$ per ogni $\omega$ tale che
    \[
      0 \le P(\omega) \le 1 \qquad\text{e}\qquad \sum_{\omega\in\Omega} P(\omega) = 1 .
    \]
  \end{itemize}
\end{defbox}

Un \textbf{evento} $A$ è un qualsiasi sottoinsieme di $\Omega$; la sua probabilità è la somma delle probabilità dei punti che contiene:
\[
  P(A) = \sum_{\omega \in A} P(\omega)
\]

\subsubsection*{Variabili aleatorie}

Una \textbf{variabile aleatoria} (VA) è una variabile il cui valore dipende dal caso: l'agente non può sceglierlo, può solo averne una credenza. Ogni VA ha un \textbf{dominio}, cioè l'insieme dei valori che può assumere.

La probabilità sui punti campionari \emph{induce} una distribuzione su ogni VA $X$: la probabilità che $X$ assuma il valore $x_i$ è la somma dei punti in cui ciò accade:
\[
  P(X = x_i) = \sum_{\omega \,:\, X(\omega) = x_i} P(\omega)
\]

\subsubsection*{Proposizioni come eventi}

Una \textbf{proposizione} si interpreta come l'evento (l'insieme dei punti campionari) in cui essa è vera. Date due VA booleane $A$ e $B$:
\begin{itemize}
  \item l'evento $a$ è l'insieme dei punti in cui $A(\omega) = \text{true}$;
  \item l'evento $\lnot a$ è l'insieme dei punti in cui $A(\omega) = \text{false}$;
  \item l'evento $a \land b$ è l'insieme dei punti in cui $A$ e $B$ sono entrambe vere.
\end{itemize}

Con variabili booleane, \textbf{un punto campionario coincide con un modello della logica proposizionale}. Di conseguenza ogni proposizione equivale alla disgiunzione degli eventi atomici in cui è vera, e la sua probabilità è la somma delle loro probabilità:
\[
  (a \lor b) \equiv (\lnot a \land b) \lor (a \land \lnot b) \lor (a \land b)
  \;\Longrightarrow\;
  P(a \lor b) = P(\lnot a \land b) + P(a \land \lnot b) + P(a \land b)
\]

\subsubsection*{Perché proprio la probabilità?}

Le definizioni impongono che eventi \emph{logicamente collegati} abbiano probabilità \emph{coerenti}; ad esempio:
\[
  P(a \lor b) = P(a) + P(b) - P(a \land b)
\]

\begin{chiave}{Argomento di de Finetti (1931)}
  Un agente che scommette secondo credenze che \textbf{violano} gli assiomi della probabilità può essere costretto a una serie di scommesse in cui \textbf{perde denaro qualunque sia l'esito}. Gli assiomi sono quindi le uniche regole razionali per gestire credenze incerte.
\end{chiave}

\subsection{Sintassi delle proposizioni}

Le proposizioni di base sono le variabili aleatorie; le proposizioni complesse sono combinazioni booleane arbitrarie di VA. Si distinguono tre tipi di VA:

\begin{center}
  \renewcommand{\arraystretch}{1.35}
  \begin{tabular}{>{\bfseries}l p{10.6cm}}
    \toprule
    \intest{Tipo} & \intest{Caratteristiche} \\
    \midrule
    Booleana & dominio $\{\text{true}, \text{false}\}$; la proposizione $X = \text{true}$ si scrive in breve $x$ (minuscolo) \\
    Discreta & dominio finito o infinito numerabile; i valori devono essere \textbf{esaustivi} e \textbf{mutuamente esclusivi} \\
    Continua & valori reali, limitati o illimitati; si ammettono anche proposizioni come $X < 22{,}0$ \\
    \bottomrule
  \end{tabular}
\end{center}

\subsubsection*{Eventi atomici}

Un \textbf{evento atomico} è un'assegnazione di valori a \emph{tutte} le variabili del dominio. Gli eventi atomici sono l'analogo probabilistico dei modelli della logica.

\begin{proprietà}{Proprietà degli eventi atomici}
  \begin{enumerate}
    \item sono \textbf{mutuamente esclusivi}: due eventi atomici diversi non possono essere veri contemporaneamente;
    \item sono \textbf{esaustivi}: la disgiunzione di tutti gli eventi atomici è sempre vera;
    \item ogni evento atomico \textbf{determina la verità di ogni proposizione} (con la semantica standard dei connettivi);
    \item ogni proposizione è logicamente equivalente alla \textbf{disgiunzione degli eventi atomici} in cui è vera.
  \end{enumerate}
\end{proprietà}

\subsection{Probabilità a priori e distribuzione congiunta}

\begin{defbox}{Probabilità a priori}
  La probabilità \textbf{a priori} (o incondizionata) $P(a)$ è la credenza in $a$ \textbf{prima} dell'arrivo di qualsiasi nuova evidenza. Svolge il ruolo dei \emph{fatti} in una KB.

  La \textbf{distribuzione di probabilità} $\Pv(X)$ (in grassetto) elenca le probabilità di tutti i valori di $X$, ed è \emph{normalizzata}: la somma vale 1. Ad esempio, per una VA con quattro valori, $\Pv(X) = \tup{0.72,\, 0.1,\, 0.08,\, 0.1}$.
\end{defbox}

\begin{defbox}{Distribuzione congiunta}
  La \textbf{distribuzione di probabilità congiunta} di un insieme di VA assegna una probabilità a \textbf{ogni evento atomico} su quelle variabili. Per due VA con 4 e 2 valori, $\Pv(X, Y)$ è una tabella $4 \times 2$.
\end{defbox}

\begin{chiave}{Potere della congiunta completa}
  La distribuzione congiunta \textbf{completa} (su tutte le variabili del dominio) permette di rispondere a \textbf{qualsiasi domanda} sul dominio, perché ogni evento è una somma di punti campionari.
\end{chiave}

\subsubsection*{Variabili continue}

Per una VA continua non si elencano i valori (sono infiniti), ma si esprime la distribuzione come una \textbf{funzione parametrica} del valore: una \textbf{densità di probabilità}, il cui integrale su tutto il dominio vale 1.

\begin{itemize}
  \item \textbf{Densità uniforme} $U[a,b](x)$: costante e pari a $\frac{1}{b-a}$ su $[a,b]$, nulla altrove.
  \item Il valore della densità in un punto \textbf{non è una probabilità} (la probabilità di un singolo punto è zero), ma un limite:
  \[
    P(X = x) \;\text{ significa }\; \lim_{dx\to 0} \frac{P(x \le X \le x + dx)}{dx}
  \]
\end{itemize}

\begin{defbox}{Densità gaussiana}
  \[
    P(x) = \frac{1}{\sqrt{2\pi}\,\sigma}\; e^{-(x-\mu)^2 / 2\sigma^2}
  \]
  con media $\mu$ e deviazione standard $\sigma$. L'area sotto la curva:
  \begin{itemize}
    \item tra $\mu-\sigma$ e $\mu+\sigma$ è circa il $68{,}2\%$;
    \item tra $\mu-2\sigma$ e $\mu+2\sigma$ è circa il $95{,}4\%$;
    \item tra $\mu-3\sigma$ e $\mu+3\sigma$ è circa il $99{,}7\%$.
  \end{itemize}
\end{defbox}

\subsection{Probabilità condizionata}

La probabilità \textbf{condizionata} (o a posteriori) $P(a \mid b)$ è la credenza in $a$ \textbf{sapendo che $b$ è tutto ciò che si conosce}.

\textbf{Attenzione}: $P(a \mid b) = 0.6$ \emph{non} significa ``se $b$, allora $a$ con probabilità 60\%''. Non è una regola che vale sempre quando $b$ è vera: vale solo quando $b$ è \emph{l'unica} informazione disponibile.

Alcune conseguenze:
\begin{itemize}
  \item se si viene a sapere di più (ad esempio che anche $a$ è vera), la credenza cambia: $P(a \mid b, a) = 1$;
  \item la credenza meno specifica \emph{resta valida} dopo nuove evidenze, ma può non essere più utile;
  \item una nuova evidenza può essere \textbf{irrilevante}, e allora la si può ignorare: $P(a \mid b, c) = P(a \mid b)$. Questo tipo di semplificazione, giustificata dalla conoscenza del dominio, è \textbf{cruciale} per rendere l'inferenza trattabile.
\end{itemize}

Notazione: $\Pv(X \mid Y)$ indica l'insieme delle distribuzioni di $X$, una per ogni valore di $Y$.

\begin{chiave}{Definizione, regola del prodotto, regola della catena}
  \textbf{Definizione}:
  \[
    P(a \mid b) = \frac{P(a \land b)}{P(b)} \qquad \text{se } P(b) \neq 0
  \]
  \textbf{Regola del prodotto} (la stessa cosa, scritta diversamente):
  \[
    P(a \land b) = P(a \mid b)\,P(b) = P(b \mid a)\,P(a)
  \]
  \textbf{Regola della catena} (applicazione ripetuta della regola del prodotto):
  \[
    \Pv(X_1, \dots, X_n) = \Pv(X_1, \dots, X_{n-1})\,\Pv(X_n \mid X_1, \dots, X_{n-1})
    = \dots = \prod_{i=1}^{n} \Pv(X_i \mid X_1, \dots, X_{i-1})
  \]
\end{chiave}

La regola del prodotto vale anche per intere distribuzioni: $\Pv(X, Y) = \Pv(X \mid Y)\,\Pv(Y)$ va letta come un \emph{insieme di equazioni}, una per ogni coppia di valori, e non come un prodotto di matrici.

\subsection{Inferenza per enumerazione}

L'inferenza probabilistica consiste nel calcolare probabilità a posteriori a partire dalle evidenze osservate. Il metodo più diretto parte dalla distribuzione congiunta completa.

\textbf{Probabilità di una proposizione.} Per ogni proposizione $\varphi$ si sommano gli eventi atomici in cui è vera:
\[
  P(\varphi) = \sum_{\omega \,:\, \omega \models \varphi} P(\omega)
\]
Il procedimento è corretto perché ogni proposizione equivale alla disgiunzione degli eventi atomici in cui vale, e questi sono mutuamente esclusivi (quindi non ci sono sovrapposizioni da contare due volte).

\textbf{Probabilità condizionate.} Si applica la definizione, calcolando numeratore e denominatore come somme di eventi atomici:
\[
  P(a \mid b) = \frac{P(a \land b)}{P(b)}
\]

\subsubsection*{Normalizzazione}

Il denominatore $P(b)$ è lo stesso per tutti i valori della variabile di interrogazione: lo si può quindi trattare come una \textbf{costante di normalizzazione} $\alpha = 1/P(b)$. Si calcolano i valori \emph{non normalizzati} e poi li si scala in modo che sommino a 1:
\[
  \Pv(X \mid b) = \alpha\, \Pv(X, b)
\]
Ad esempio, se si ottiene $\alpha\tup{0.12,\, 0.08}$, normalizzando si ha $\tup{0.6,\, 0.4}$. Il vantaggio è che non serve calcolare $P(b)$ esplicitamente.

\subsubsection*{Formula generale}

Si suddividono tutte le variabili $\mathbf{X}$ in tre gruppi:
\begin{itemize}
  \item $\mathbf{Y}$: variabili di \textbf{interrogazione} (ciò che si vuole sapere);
  \item $\mathbf{E}$: variabili di \textbf{evidenza}, con valori osservati $\mathbf{e}$;
  \item $\mathbf{H} = \mathbf{X} - \mathbf{Y} - \mathbf{E}$: variabili \textbf{nascoste} (né chieste né osservate).
\end{itemize}

\begin{chiave}{Inferenza per enumerazione}
  \[
    \Pv(\mathbf{Y} \mid \mathbf{E} = \mathbf{e})
    = \alpha\, \Pv(\mathbf{Y}, \mathbf{E} = \mathbf{e})
    = \alpha \sum_{\mathbf{h}} \Pv(\mathbf{Y}, \mathbf{E} = \mathbf{e}, \mathbf{H} = \mathbf{h})
  \]
  Si fissano le evidenze e si \textbf{sommano via} (marginalizzano) le variabili nascoste. I termini della somma sono voci della congiunta, perché $\mathbf{Y}$, $\mathbf{E}$ e $\mathbf{H}$ insieme esauriscono tutte le variabili.
\end{chiave}

\textbf{Limiti evidenti}, con $n$ variabili e $d$ valori al massimo per variabile:
\begin{enumerate}
  \item complessità temporale nel caso peggiore $O(d^n)$;
  \item complessità spaziale $O(d^n)$ per memorizzare la congiunta;
  \item soprattutto: \textbf{da dove si ricavano} $O(d^n)$ numeri?
\end{enumerate}
La soluzione è sfruttare la struttura del dominio, cioè le \emph{indipendenze}.

\subsection{Indipendenza}

\begin{defbox}{Indipendenza (assoluta)}
  $A$ e $B$ sono \textbf{indipendenti} se e solo se vale una delle seguenti condizioni (equivalenti):
  \[
    \Pv(A \mid B) = \Pv(A) \qquad \Pv(B \mid A) = \Pv(B) \qquad \Pv(A, B) = \Pv(A)\,\Pv(B)
  \]
  Intuitivamente: sapere qualcosa su $B$ non cambia la credenza su $A$.
\end{defbox}

L'indipendenza permette di \textbf{fattorizzare} la congiunta in tabelle più piccole. Se un gruppo di 3 variabili booleane è indipendente da una variabile con 4 valori, la congiunta passa da $2^3 \cdot 4 = 32$ voci a $2^3 + 4 = 12$. Per $n$ variabili booleane tutte indipendenti si passa da $2^n$ a $n$ numeri.

Tuttavia l'indipendenza assoluta è \textbf{potente ma rara}: nei domini reali le variabili sono quasi sempre collegate tra loro in qualche modo.

\subsection{Indipendenza condizionata}

Spesso due variabili non sono indipendenti in assoluto, ma lo diventano \textbf{una volta nota la loro causa comune}. Se due effetti dipendono entrambi da una stessa causa, una volta conosciuta la causa sapere un effetto non dice nulla di nuovo sull'altro.

\begin{chiave}{Indipendenza condizionata}
  $X$ è \textbf{condizionatamente indipendente} da $Y$ dato $Z$ se:
  \[
    \Pv(X \mid Y, Z) = \Pv(X \mid Z)
  \]
  Formulazioni equivalenti:
  \[
    \Pv(Y \mid X, Z) = \Pv(Y \mid Z)
    \qquad\qquad
    \Pv(X, Y \mid Z) = \Pv(X \mid Z)\,\Pv(Y \mid Z)
  \]
\end{chiave}

\textbf{Effetto sulla rappresentazione.} Con tre variabili booleane (una causa $C$ e due effetti $E_1, E_2$), la congiunta ha $2^3 - 1 = 7$ numeri indipendenti (l'ottavo si ricava perché la somma vale 1). Applicando la regola della catena e poi l'indipendenza condizionata:
\begin{align*}
  \Pv(E_1, E_2, C)
  &= \Pv(E_1 \mid E_2, C)\,\Pv(E_2, C) \\
  &= \Pv(E_1 \mid E_2, C)\,\Pv(E_2 \mid C)\,\Pv(C) \\
  &= \Pv(E_1 \mid C)\,\Pv(E_2 \mid C)\,\Pv(C)
\end{align*}
Servono ora $2 + 2 + 1 = 5$ numeri invece di 7.

In generale l'indipendenza condizionata riduce la dimensione della rappresentazione della congiunta da \textbf{esponenziale} in $n$ a \textbf{lineare} in $n$. È la forma di conoscenza più basilare e robusta sugli ambienti incerti.

\subsection{Regola di Bayes}

Dalla regola del prodotto, $P(a \land b) = P(a \mid b)P(b) = P(b \mid a)P(a)$; uguagliando i due membri e dividendo per $P(b)$ si ottiene:

\begin{chiave}{Regola di Bayes}
  \[
    P(a \mid b) = \frac{P(b \mid a)\,P(a)}{P(b)}
    \qquad\qquad
    \Pv(Y \mid X) = \frac{\Pv(X \mid Y)\,\Pv(Y)}{\Pv(X)} = \alpha\, \Pv(X \mid Y)\,\Pv(Y)
  \]
\end{chiave}

\textbf{Perché è utile.} Nella pratica si conosce più facilmente la probabilità \emph{causale} (quanto spesso una causa produce un effetto), mentre serve quella \emph{diagnostica} (quanto è probabile la causa, visto l'effetto). La regola di Bayes permette di passare dall'una all'altra:
\[
  P(\mathit{Causa} \mid \mathit{Effetto}) = \frac{P(\mathit{Effetto} \mid \mathit{Causa})\,P(\mathit{Causa})}{P(\mathit{Effetto})}
\]
Si noti il ruolo della probabilità a priori $P(\mathit{Causa})$: se la causa è molto rara, la probabilità a posteriori può restare piccola anche quando l'effetto è fortemente associato alla causa.

\subsubsection*{Bayes e indipendenza condizionata: il modello naive Bayes}

Combinando la regola di Bayes con l'indipendenza condizionata degli effetti data la causa:
\[
  \Pv(C \mid e_1 \land e_2) = \alpha\, \Pv(e_1 \land e_2 \mid C)\,\Pv(C) = \alpha\, \Pv(e_1 \mid C)\,\Pv(e_2 \mid C)\,\Pv(C)
\]

\begin{defbox}{Modello naive Bayes}
  Una causa influenza direttamente più effetti, che sono condizionatamente indipendenti tra loro data la causa:
  \[
    \Pv(\mathit{Causa}, \mathit{Effetto}_1, \dots, \mathit{Effetto}_n) = \Pv(\mathit{Causa}) \prod_{i} \Pv(\mathit{Effetto}_i \mid \mathit{Causa})
  \]
  Il numero di parametri cresce \textbf{linearmente} con il numero di effetti. Si dice ``naive'' perché spesso viene usato anche quando l'indipendenza non è del tutto verificata, ottenendo comunque buoni risultati.
\end{defbox}

\clearpage
\subsection{Il mondo del Wumpus con l'incertezza}

Il Wumpus è l'esempio che unisce tutti gli strumenti visti: mostra dove la logica si blocca e come la probabilità, insieme all'indipendenza condizionata, permette comunque di decidere.

\subsubsection*{Descrizione PEAS}

\begin{center}
  \renewcommand{\arraystretch}{1.35}
  \begin{tabular}{>{\bfseries}l p{11.2cm}}
    \toprule
    \intest{Componente} & \intest{Descrizione} \\
    \midrule
    Prestazione & oro $+1000$, morte $-1000$; $-1$ per ogni passo, $-10$ per l'uso della freccia \\
    Ambiente & griglia $4\times4$. Le caselle adiacenti al Wumpus hanno \emph{puzza}; quelle adiacenti a un pozzo hanno \emph{brezza}; c'è \emph{luccichio} se e solo se l'oro è nella stessa casella. Sparare uccide il Wumpus se l'agente lo ha di fronte; c'è una sola freccia. \emph{Grab} raccoglie l'oro nella stessa casella, \emph{Release} lo lascia \\
    Attuatori & gira a sinistra, gira a destra, avanti, \emph{Grab}, \emph{Release}, \emph{Shoot} \\
    Sensori & brezza, luccichio, puzza \\
    \bottomrule
  \end{tabular}
\end{center}

\subsubsection*{Una situazione senza via d'uscita logica}

L'agente ha visitato $[1,1]$, $[1,2]$ e $[2,1]$: nessuna contiene un pozzo (è ancora vivo), in $[1,1]$ non c'è brezza, mentre c'è brezza sia in $[1,2]$ sia in $[2,1]$.

Le caselle di frontiera sono $[1,3]$, $[2,2]$ e $[3,1]$. La logica dice soltanto che c'è un pozzo in $[1,3]$ o in $[2,2]$ (per la brezza in $[1,2]$) e un pozzo in $[2,2]$ o in $[3,1]$ (per la brezza in $[2,1]$). \textbf{Nessuna casella è dimostrabilmente sicura}: un agente puramente logico non ha azioni sicure e non sa scegliere.

Eppure le caselle non sono tutte equivalenti: il calcolo che segue mostra che, con pozzi distribuiti uniformemente, $[2,2]$ contiene un pozzo con probabilità $\approx 0.86$, mentre $[1,3]$ e $[3,1]$ solo con probabilità $\approx 0.31$. L'agente razionale (MEU) andrà quindi in $[1,3]$ o in $[3,1]$.

\subsubsection*{Variabili del modello}

\begin{itemize}
  \item $P_{i,j} = \text{true}$ se e solo se la casella $[i,j]$ contiene un pozzo (16 variabili);
  \item $B_{i,j} = \text{true}$ se e solo se la casella $[i,j]$ ha brezza; nel modello si includono soltanto quelle osservate, $B_{1,1}$, $B_{1,2}$, $B_{2,1}$.
\end{itemize}

\subsubsection*{Specifica del modello probabilistico}

La congiunta completa è $\Pv(P_{1,1}, \dots, P_{4,4}, B_{1,1}, B_{1,2}, B_{2,1})$. Invece di assegnarla direttamente, la si scompone con la regola del prodotto, in modo da ottenere la forma $P(\mathit{Effetto} \mid \mathit{Causa})$, che è facile da specificare:
\[
  \Pv(B_{1,1}, B_{1,2}, B_{2,1} \mid P_{1,1}, \dots, P_{4,4})\;\Pv(P_{1,1}, \dots, P_{4,4})
\]

\begin{itemize}
  \item \textbf{Primo fattore} (effetto dato la causa): è \textbf{deterministico}. Vale 1 se le brezze sono coerenti con i pozzi (c'è brezza esattamente nelle caselle adiacenti a qualche pozzo), 0 altrimenti.
  \item \textbf{Secondo fattore} (a priori sui pozzi): ogni casella contiene un pozzo con probabilità $0.2$, \textbf{indipendentemente} dalle altre. Per una configurazione con $n$ pozzi:
  \[
    \Pv(P_{1,1}, \dots, P_{4,4}) = \prod_{i,j=1,1}^{4,4} P(P_{i,j}) = 0.2^{\,n} \times 0.8^{\,16-n}
  \]
\end{itemize}

\subsubsection*{Osservazioni e interrogazione}

I fatti noti sono:
\[
  b = \lnot b_{1,1} \land b_{1,2} \land b_{2,1}
  \qquad\qquad
  \known = \lnot p_{1,1} \land \lnot p_{1,2} \land \lnot p_{2,1}
\]
L'interrogazione è $\Pv(P_{1,3} \mid \known, b)$. Si chiamano $\mathit{Unknown}$ tutte le variabili $P_{i,j}$ diverse da $P_{1,3}$ e da quelle in $\mathit{Known}$ (sono 12). Per enumerazione:
\[
  \Pv(P_{1,3} \mid \known, b) = \alpha \sum_{\unknown} \Pv(P_{1,3}, \unknown, \known, b)
\]
Il problema è che la somma ha $2^{12} = 4096$ termini: \textbf{cresce esponenzialmente con il numero di caselle}.

\subsubsection*{L'idea chiave: indipendenza condizionata}

Le brezze osservate dipendono \emph{soltanto} dalle caselle vicine. Si divide perciò $\mathit{Unknown}$ in due parti:
\begin{itemize}
  \item $\mathit{Fringe}$ (frontiera): le caselle sconosciute adiacenti a quelle visitate, escluse quella interrogata; qui $\{P_{2,2}, P_{3,1}\}$;
  \item $\mathit{Other}$: tutte le altre caselle sconosciute (10 variabili).
\end{itemize}

\begin{chiave}{Ipotesi di indipendenza condizionata}
  Le osservazioni sono condizionatamente indipendenti dalle caselle lontane, date le caselle vicine:
  \[
    \Pv(b \mid P_{1,3}, \mathit{Known}, \mathit{Unknown}) = \Pv(b \mid P_{1,3}, \mathit{Known}, \mathit{Fringe})
  \]
\end{chiave}

\subsubsection*{La derivazione passo per passo}

Si manipola l'interrogazione fino a poter sfruttare questa ipotesi. Ogni passaggio è numerato e giustificato sotto.
{\small
\begin{align}
  \Pv(P_{1,3} \mid \known, b)
  &= \alpha \sum_{\unknown} \Pv(P_{1,3}, \unknown, \known, b) \\
  &= \alpha \sum_{\unknown} \Pv(b \mid P_{1,3}, \known, \unknown)\, \Pv(P_{1,3}, \known, \unknown) \\
  &= \alpha \sum_{\fringe} \sum_{\other} \Pv(b \mid \known, P_{1,3}, \fringe, \other)\, \Pv(P_{1,3}, \known, \fringe, \other) \\
  &= \alpha \sum_{\fringe} \sum_{\other} \Pv(b \mid \known, P_{1,3}, \fringe)\, \Pv(P_{1,3}, \known, \fringe, \other) \\
  &= \alpha \sum_{\fringe} \Pv(b \mid \known, P_{1,3}, \fringe) \sum_{\other} \Pv(P_{1,3}, \known, \fringe, \other) \\
  &= \alpha \sum_{\fringe} \Pv(b \mid \known, P_{1,3}, \fringe) \sum_{\other} \Pv(P_{1,3})\,P(\known)\,P(\fringe)\,P(\other) \\
  &= \alpha\, P(\known)\,\Pv(P_{1,3}) \sum_{\fringe} \Pv(b \mid \known, P_{1,3}, \fringe)\,P(\fringe) \sum_{\other} P(\other) \\
  &= \alpha'\, \Pv(P_{1,3}) \sum_{\fringe} \Pv(b \mid \known, P_{1,3}, \fringe)\,P(\fringe)
\end{align}
}

\begin{enumerate}
  \item \textbf{enumerazione}: si sommano via le variabili nascoste;
  \item \textbf{regola del prodotto}, per isolare il termine $P(\mathit{Effetto} \mid \mathit{Causa})$;
  \item si scompone $\mathit{Unknown} = \mathit{Fringe} \cup \mathit{Other}$ in due somme;
  \item \textbf{indipendenza condizionata}: $b$ non dipende da $\mathit{Other}$, data la frontiera;
  \item il primo fattore non dipende da $\other$ e si porta fuori dalla somma interna;
  \item i pozzi sono \textbf{indipendenti} tra loro, quindi la loro a priori si fattorizza;
  \item si raccolgono fuori dalle somme i fattori che non dipendono dagli indici di somma;
  \item $\sum_{\other} P(\other) = 1$ e $P(\known)$ viene assorbito in $\alpha' = \alpha\,P(\known)$.
\end{enumerate}

Negli ultimi due passaggi accadono due cose:
\begin{itemize}
  \item $\sum_{\other} P(\other) = 1$, perché è la somma di una distribuzione su tutti i suoi valori: le 10 caselle lontane \textbf{spariscono} dal calcolo;
  \item $P(\known)$ è una costante (non dipende da $P_{1,3}$), quindi viene assorbita nella nuova costante di normalizzazione $\alpha' = \alpha\,P(\known)$.
\end{itemize}

Il risultato è che la somma corre soltanto sulle $2^2 = 4$ configurazioni della frontiera, non più su $2^{12}$.

\subsubsection*{Il calcolo numerico}

Il fattore $\Pv(b \mid \known, P_{1,3}, \fringe)$ vale 1 per le configurazioni coerenti con le brezze osservate e 0 per le altre. Basta quindi sommare $P(\fringe)$ sulle configurazioni \textbf{coerenti}; ogni casella di frontiera contribuisce con $0.2$ se ha un pozzo e $0.8$ se non lo ha.

Le condizioni da rispettare sono: brezza in $[1,2]$ $\Rightarrow$ pozzo in $[1,3]$ o in $[2,2]$; brezza in $[2,1]$ $\Rightarrow$ pozzo in $[2,2]$ o in $[3,1]$.

\begin{center}
  \renewcommand{\arraystretch}{1.3}
  \begin{tabular}{c c c c c}
    \toprule
    \intest{$P_{1,3}$} & \intest{$P_{2,2}$} & \intest{$P_{3,1}$} & \intest{Coerente?} & \intest{$P(\fringe)$} \\
    \midrule
    \vero & \vero & \vero & sì & $0.2 \cdot 0.2 = 0.04$ \\
    \vero & \vero & \falso & sì & $0.2 \cdot 0.8 = 0.16$ \\
    \vero & \falso & \vero & sì & $0.8 \cdot 0.2 = 0.16$ \\
    \vero & \falso & \falso & \nota{no: $[2,1]$ senza causa} & -- \\
    \midrule
    \falso & \vero & \vero & sì & $0.04$ \\
    \falso & \vero & \falso & sì & $0.16$ \\
    \falso & \falso & \vero & \nota{no: $[1,2]$ senza causa} & -- \\
    \falso & \falso & \falso & \nota{no} & -- \\
    \bottomrule
  \end{tabular}
\end{center}

Moltiplicando per la probabilità a priori di $P_{1,3}$ ($0.2$ se vero, $0.8$ se falso) e normalizzando:

\begin{chiave}{Risultato}
  \[
    \Pv(P_{1,3} \mid \known, b) = \alpha' \tup{0.2\,(0.04 + 0.16 + 0.16),\; 0.8\,(0.04 + 0.16)}
    = \alpha' \tup{0.072,\; 0.16} \approx \tup{0.31,\; 0.69}
  \]
  Per simmetria lo stesso vale per $[3,1]$. Ripetendo il calcolo per $[2,2]$ (frontiera $\{P_{1,3}, P_{3,1}\}$):
  \[
    \Pv(P_{2,2} \mid \known, b) = \alpha' \tup{0.2 \cdot 1,\; 0.8 \cdot 0.04} = \alpha'\tup{0.2,\; 0.032} \approx \tup{0.86,\; 0.14}
  \]
\end{chiave}

\textbf{Interpretazione.} Se $[2,2]$ ha un pozzo, entrambe le brezze sono spiegate da sole, qualunque sia il contenuto delle altre caselle (somma $=1$). Se invece $[2,2]$ è libera, servono \emph{contemporaneamente} due pozzi, in $[1,3]$ e in $[3,1]$: un evento poco probabile ($0.04$). Per questo $[2,2]$ è la casella più pericolosa, e l'agente dovrebbe scegliere $[1,3]$ o $[3,1]$.

\subsection{Riepilogo}

\begin{riepilogo}{Punti fondamentali}
  \begin{itemize}
    \item La probabilità è un formalismo rigoroso per rappresentare conoscenza incerta.
    \item La distribuzione congiunta assegna una probabilità a ogni evento atomico.
    \item Le interrogazioni si risolvono sommando su eventi atomici (inferenza per enumerazione).
    \item Nei domini non banali la congiunta è troppo grande: bisogna ridurne la dimensione.
    \item L'indipendenza e, soprattutto, l'indipendenza condizionata forniscono gli strumenti per farlo.
  \end{itemize}
\end{riepilogo}

\end{document}
